\begin{claim}{C.20}{Refuted}{Pure-injective tests do not suffice for arbitrary sums of representable components}\label{C20}
There exist a commutative ring $R$, a projective $R$-module $Q$, and an exact $R$-linear functor
\[
 E:\Mod R\longrightarrow\Ch_{\geq0}(R)
\]
such that every $E_j$ is a countable direct sum of functors $\operatorname{Hom}_R(P,-)$ with $P$ projective, every $E(I)$ is acyclic for pure-injective $I$, but $H_0E(Q)\ne0$. Consequently the finite representability proof of Claim~\ref{C16} cannot be extended to arbitrary sums using only exactness, projective representing modules, and pure-injective tests. This is a counterexample to that proposed proof principle; we do not identify $E$ with an adjunction satisfying all the manuscript's monoidal hypotheses.
\end{claim}
\begin{proof}
We first specify the coefficient modules and the finite resolutions that will be reversed. For $n\geq1$, let
\[
 R_n=\mathbb F_2[x_\alpha\mid\alpha<\aleph_n]
       /(x_\alpha^2-x_\alpha\mid\alpha<\aleph_n),
 \qquad L_n=R_n/(x_\alpha\mid\alpha<\aleph_n).
\]
Thus $R_n$ is the unital free Boolean ring on $\aleph_n$ generators and $L_n\cong\mathbb F_2$ is its augmentation residue. We use two precise classical dimension statements. The free Boolean ring on $\aleph_n$ generators has global dimension $n+1$, as stated in \citet[p.~319, following Corollary 1.5]{Osofsky1968C}, with attribution there to Pierce. The ideal generated by an independent family of $\aleph_n$ idempotents has projective dimension at least $n$, as stated in \citet[Remark 5.1(e), p.~917]{FinnMartinezMcGovern1997C}, with attribution there to Osofsky. Here independence means that, for disjoint finite subsets $A,B$ of the family, the product $\prod_{a\in A}a\prod_{b\in B}(1-b)$ is nonzero; this is precisely Remark 5.1(c) of that opened source. The free generators form such an independent family, as one sees by evaluating the variables in $A$ at $1$ and those in $B$ at $0$. The short exact sequence
\[
 0\longrightarrow(x_\alpha)\longrightarrow R_n\longrightarrow L_n\longrightarrow0
\]
therefore gives $\operatorname{pd}_{R_n}L_n=n+1$: the upper bound is the global-dimension statement; a smaller bound would give projective dimension at most $n-1$ for the ideal, contradicting the second statement. We use these opened formulations of the classical estimates; the original Pierce article and Osofsky monograph were not obtained and are not reported as opened proofs.

Put $r_n=n+1$. Take the standard free resolution of $L_n$ obtained successively by taking the free module on the underlying set of each kernel. Since its projective dimension is $r_n$, the kernel in position $r_n$ is projective. Stopping with that kernel, rather than deleting the top of an infinite resolution, gives a finite exact sequence
\[
 0\longrightarrow P^{(n)}_{r_n}\xrightarrow{d_{r_n}}
 P^{(n)}_{r_n-1}\longrightarrow\cdots\longrightarrow
 P^{(n)}_0\longrightarrow L_n\longrightarrow0,
 \tag{C-finite-resolution}\label{Cfiniteresolution}
\]
with every term projective. The identity of $P^{(n)}_{r_n}$ determines a nonzero class
\[
 \eta_n\in\operatorname{Ext}^{r_n}_{R_n}
             (L_n,P^{(n)}_{r_n}).
\]
Indeed a zero class would mean that the identity extends across $d_{r_n}$, splitting the top inclusion and shortening the resolution, contrary to $\operatorname{pd}L_n=r_n$.

Set $R=\prod_{n\geq1}R_n$ and let $e_n\in R$ be the coordinate idempotents. We regard all modules in \eqref{Cfiniteresolution} as $R$-modules supported on $e_n$. They remain projective whenever they were projective over $R_n$, because $e_nR\cong R_n$ is a projective $R$-module. Define
\[
 Q=\bigoplus_{n\geq1}P^{(n)}_{r_n},
\]
a projective $R$-module. The displayed classes remain nonzero in $\operatorname{Ext}^{r_n}_R(L_n,Q)$, since applying Hom from the $e_n$-supported resolution sees precisely $e_nQ=P^{(n)}_{r_n}$.

The ring $R$ is Boolean, so it is reduced and all its prime quotients are $\mathbb F_2$. It is absolutely flat by \citet[Lemma 15.106.5, tag 092F]{StacksAbsolutelyFlatC}. Thus every short exact sequence of $R$-modules is pure, and pure-injective modules are exactly injective modules.

For $X\in\Mod R$ define a finite nonnegative chain complex by reversing the entire Hom complex of \eqref{Cfiniteresolution}:
\[
 C_n(X)_j=\operatorname{Hom}_R(P^{(n)}_{r_n-j},X)
 \quad(0\leq j\leq r_n),
\]
and set the other terms to zero. Its differential is precomposition with the resolution differential. Then
\[
 H_jC_n(X)=\operatorname{Ext}^{r_n-j}_R(L_n,X)
 \quad(0\leq j\leq r_n).
\]
In particular, for injective $I$ its homology is concentrated in degree $r_n$, where it is $\operatorname{Hom}_R(L_n,I)$. The finite top projective term in \eqref{Cfiniteresolution} is essential for this assertion: a brutal truncation of an infinite resolution would introduce a spurious degree-zero cokernel.

Let $A(X)=\bigoplus_{n\geq1}C_n(X)$, $B(X)=\prod_{n\geq1}C_n(X)$, and let $\iota_X:A(X)\to B(X)$ be the natural inclusion. Define
\[
 E(X)=\operatorname{Cone}(\iota_X),\qquad
 E(X)_j=B(X)_j\oplus A(X)_{j-1}.
\]
The differential is $(b,a)\mapsto(d_Bb+\iota_Xa,-d_Aa)$. Every component of $E$ is exact, because all the $P^{(n)}_i$ are projective and sums and products of module exact sequences are exact. More precisely,
\[
 E_j\cong
 \operatorname{Hom}_R\!\left(\bigoplus_nP^{(n)}_{r_n-j},-\right)
 \ \oplus\ \bigoplus_n
 \operatorname{Hom}_R(P^{(n)}_{r_n-j+1},-),
\]
where an out-of-range resolution term is zero. This is a countable sum of representables with projective representing modules, as asserted.

For injective $I$, homology commutes with both the indicated sum and product. For each fixed degree $j$, at most one $n$ has $r_n=j$. Hence $H_j(\iota_I)$ is an isomorphism: the sum and product of the homology groups both have at most one nonzero contribution. Its cone $E(I)$ is acyclic. This covers every pure-injective $I$ over $R$.

For the projective module $Q$, connectivity and the cone homology sequence instead give
\[
 H_0E(Q)\cong
 \frac{\displaystyle\prod_{n\geq1}
          \operatorname{Ext}^{r_n}_R(L_n,Q)}
      {\displaystyle\bigoplus_{n\geq1}
          \operatorname{Ext}^{r_n}_R(L_n,Q)}\ne0.
\]
The sequence $(\eta_n)_n$ represents a nonzero class because every coordinate is nonzero, whereas the denominator consists of finite-support sequences. This proves the counterexample and the stated limitation of the pro-representability route. The coefficient ring and modules here are intrinsically infinite; no finite computation is being substituted for the cardinal dimension theorems or the product/sum argument.
\end{proof}

